\section{Observation model and sensitivity assumption}\label{sec:model}
Fix a terminal time $H$ and choose times $0=H_0<H_1<\ldots<H_K=H$. Let $L$ take value 1 if an adverse outcome occurs and 0 otherwise. Let $J_\Pi=k$ denote the event that verification terminates within $(H_{k-1},H_k]$ for some $k\le K$, and $J_\Pi=\infty$ if it does not terminate before $H$. The observed data value is $Y=(k,L)$ if $J_\Pi\le K$, or $Y=\infty$ if $J_\Pi=\infty$. We therefore have $d=2K+1$ possible values for the observed distribution of $Y$:
\begin{equation}\label{eq:masses}
 a_k=\Pbb(J_\Pi=k,L=1),\qquad b_k=\Pbb(J_\Pi=k,L=0),\qquad c=\Pbb(J_\Pi=\infty).
\end{equation}
Set $A=\sum_k a_k$, $G=\sum_k b_k$, and $r=A+G=1-c$. The only unidentified cell split is
\begin{equation}\label{eq:hidden}
 u=\Pbb(J_\Pi=\infty,L=1)\in[0,c],\qquad \theta=A+u.
\end{equation}
Labels considered resolved will be taken as ground truth. We specify the set of rules for horizon, partitioning and observations ahead of time prior to knowledge of these outcomes.

We similarly pre-specify the finest granularity $J^*$ at which we represent categories. For any given partitioning $\Pi$, there is an associated mapping $g_\Pi: J^* \mapsto J_\Pi$ resulting in a coarsened version of $J^*$ with the unresolved label and end horizon retained. We assume: 
\begin{equation}\label{eq:pis}
 I(L;J^*)\le\rho,\qquad \rho\ge0.
\end{equation}
All logs are taken wrt $e$. For any given $\rho$, each coarsening $J_\Pi$ satisfies $I(L;J_\Pi) \leq \rho$. The sharp set of distributions of $(L, J_\Pi)$ with the observed margins on $\Pi$ and obeying this restriction will be called the sharp interval. If additional, finer margins are observed, the resulting set of compatible distributions may be reduced. We therefore view a coarse interval as corresponding to an experiment with less information than that finer one.

We use the continuous binary entropy $h(x)=-x\log x-(1-x)\log(1-x)$, with $0\log0=0$, and the homogeneous mass entropy
\begin{equation}\label{eq:massentropy}
 q(a,b)=(a+b)h\{a/(a+b)\},\qquad q(0,0)=0.
\end{equation}
Let $T_\Pi=\sum_k q(a_k,b_k)$. As before, we leave empty bands' corresponding fractions undefined since they take value 0. There is then a separate budget on compatible risk $\beta \in [0,1]$, which differs from $\rho$, and our population claim will hold if compatible risks exist and the maximum compatible risk is $\le \beta$.

$\rho$ is a sensitivity parameter. Allowing it to increase accommodates stronger degrees of selection bias (and so potentially larger identified intervals); reducing $\rho$ may be inconsistent with the observed information. We might specify several appropriate budgets using independent validation (albeit subject to concerns over finite-sample error and distribution shift), but in general can only use the partial information provided by our stream. Note that low levels of absolute information do not always rule out high degrees of selection for a low-probability event ($I(L;J^*)\le h(\theta)$ is still possible). Thus these should be considered as sensitivity analysis parameters only.
